\subsection{Direct Methods}

% **Role:** quantitative matrix behavior and computation.
%
% Direct solves, QR algorithms, SVD numerics, least-squares stability, sparse/banded structure, Krylov methods, iterative methods, preconditioning preview.

Numerical linear algebra turns matrix structure into algorithms for direct solves, least squares, eigenvalues, SVD, structured transforms, Krylov methods, iteration, and preconditioning.

\begin{perspective} \label{per:solve-with-factors-not-inverses} \hlt{Solve with Factors Before Forming Inverses}\\
For $A\mathbf{x}=\mathbf{b}$, use available factors such as LU, QR, or Cholesky to solve triangular or orthogonal subproblems.\\
This avoids forming $A^{-1}$ and preserves sparse or triangular structure.
\end{perspective}

\begin{perspective} \label{per:iterative-refinement-residual-correction} \hlt{Iterative Refinement by Residual Correction}\\
Given an approximate solution $\mathbf{x}_k$ of $A\mathbf{x}=\mathbf{b}$, set $\mathbf{r}_k=\mathbf{b}-A\mathbf{x}_k$.\\
Solve $A\mathbf{d}_k=\mathbf{r}_k$ with the same factors or preconditioner, then update $\mathbf{x}_{k+1}=\mathbf{x}_k+\mathbf{d}_k$.\\
Exact arithmetic corrects in one step; floating point helps when $\mathbf{r}_k$ is accurate and conditioning is moderate.
\end{perspective}

\begin{method} \label{meth:augmented-inverse-computation-preview} \hlt{Augmented Inverse Computation Preview}\\
For $A\in\operatorname{Mat}(n,n,\mathbb{F})$:
\begin{enumerate}[label=(\roman*)]
\setlength{\itemsep}{0pt}
\item Form $[\,A\ I_n\,]$.
\item Apply Gauss--Jordan reduction to the left block.
\item If $[\,A\ I_n\,]\sim[\,I_n\ B\,]$, then $B=A^{-1}$.
\item Otherwise $A$ is not invertible, by Corollary \ref{cor:square-matrix-invertible-iff-full-rank}.
\end{enumerate}
\end{method}

\begin{proposition} \label{prop:qr-least-squares-and-projection} \hlt{QR Route for Least Squares}\\
Let $A=QR\in\operatorname{Mat}(m,n,\mathbb{F})$ be a full-column-rank narrow QR factorisation. For $\mathbf{y}\in\mathbb{F}^m$, the least-squares solution solves $R\mathbf{x}=Q^*\mathbf{y}$. Moreover $P_{\operatorname{col}(A)}=QQ^*$.
\begin{proof}
The normal equations reduce by $A=QR$ and $Q^*Q=I_n$ to $R^*R\mathbf{x}=R^*Q^*\mathbf{y}$.
Since $R$ is invertible, this is equivalent to $R\mathbf{x}=Q^*\mathbf{y}$. Also $\operatorname{col}(A)=\operatorname{col}(Q)$, and the columns of $Q$ are orthonormal, so $P_{\operatorname{col}(A)}=QQ^*$.\\
\hlt{Structural Commentary}: QR replaces $\operatorname{col}(A)$ by an orthonormal basis $Q$.\\
\hlt{Key Insight}: Project onto $Q$, then solve triangularly.
\end{proof}
\end{proposition}

Successive orthogonalization fits in an orthogonal basis for the same column space.

\begin{breakablealgorithm}
\caption{Least Squares by Successive Orthogonalization}
\label{alg:sl-successive-orthogonalization-least-squares}
\begin{algorithmic}
\Require Full-column-rank design columns $x_0,\ldots,x_p$ and response vector $y$
\Ensure Least-squares coefficients $\widehat{\beta}$
\For{$j=0,\ldots,p$}
\State $z_j\gets x_j$
\For{$\ell=0,\ldots,j-1$}
\State $r_{\ell j}\gets q_\ell^Tx_j$
\State $z_j\gets z_j-r_{\ell j}q_\ell$
\EndFor
\State $r_{jj}\gets \norm{z_j}_2$
\State $q_j\gets z_j/r_{jj}$
\EndFor
\State Form $Q=[q_0,\ldots,q_p]$ and the upper-triangular matrix $R=(r_{\ell j})$
\State Solve $R\widehat{\beta}=Q^Ty$ by back-substitution
\end{algorithmic}
\end{breakablealgorithm}

\begin{perspective} \label{per:normal-equations-squaring-condition-number} \hlt{Why Normal Equations Can Be Numerically Fragile}\\
For full-column-rank least squares, solving $A^*A\mathbf{x}=A^*\mathbf{y}$ squares the spectral condition number because $\kappa_2(A^*A)=\kappa_2(A)^2$. QR-based methods avoid this squaring and therefore preserve conditioning better.
\end{perspective}

\begin{perspective} \label{per:linear-model-numerical-solution-route} \hlt{Solve Least Squares with Factors}\\
For $X\in\R^{n\times p}$ with $\operatorname{rank}(X)=p$, normal equations use $X^tX\widehat{\beta}=X^tY$ but $\kappa_2(X^tX)=\kappa_2(X)^2$.\\
Prefer $X=QR$ with $Q^tQ=I_p$ and $R$ nonsingular, giving $\widehat{\beta}=R^{-1}Q^tY$.\\
Use Cholesky on $X^tX$ only when $X^tX$ is well conditioned.
\end{perspective}

\begin{figure}[H]
  \centering
  \begin{tikzpicture}[
  font=\small,
  cell/.style={minimum width=0.48cm, minimum height=0.40cm, inner sep=1pt, outer sep=0pt, anchor=center},
  aband/.style={fill=gray!16},
  lband/.style={fill=blue!10, text=blue!70!black},
  uband/.style={fill=gray!16},
  title/.style={font=\small}
]
  \matrix (A) [
    matrix of math nodes,
    nodes={cell},
    left delimiter={[},
    right delimiter={]},
    row sep=0pt,
    column sep=0pt
  ] {
    |[aband]| * & |[aband]| * & \phantom{*} & \phantom{*} & \phantom{*}\\
    |[aband]| * & |[aband]| * & |[aband]| * & \phantom{*} & \phantom{*}\\
    \phantom{*} & |[aband]| * & |[aband]| * & |[aband]| * & \phantom{*}\\
    \phantom{*} & \phantom{*} & |[aband]| * & |[aband]| * & |[aband]| *\\
    \phantom{*} & \phantom{*} & \phantom{*} & |[aband]| * & |[aband]| *\\
  };
  \node[title, above=5pt] at (A.north) {$A$};

  \node at (2.35,0) {$=$};

  \begin{scope}[shift={(4.35,0)}]
  \matrix (L) [
    matrix of math nodes,
    nodes={cell},
    left delimiter={[},
    right delimiter={]},
    row sep=0pt,
    column sep=0pt
  ] {
    |[lband]| 1 & \phantom{*} & \phantom{*} & \phantom{*} & \phantom{*}\\
    |[lband]| * & |[lband]| 1 & \phantom{*} & \phantom{*} & \phantom{*}\\
    \phantom{*} & |[lband]| * & |[lband]| 1 & \phantom{*} & \phantom{*}\\
    \phantom{*} & \phantom{*} & |[lband]| * & |[lband]| 1 & \phantom{*}\\
    \phantom{*} & \phantom{*} & \phantom{*} & |[lband]| * & |[lband]| 1\\
  };
  \node[title, above=5pt] at (L.north) {$L$};
  \end{scope}

  \node at (6.55,0) {$\cdot$};

  \begin{scope}[shift={(8.55,0)}]
  \matrix (U) [
    matrix of math nodes,
    nodes={cell},
    left delimiter={[},
    right delimiter={]},
    row sep=0pt,
    column sep=0pt
  ] {
    |[uband]| * & |[uband]| * & \phantom{*} & \phantom{*} & \phantom{*}\\
    \phantom{*} & |[uband]| * & |[uband]| * & \phantom{*} & \phantom{*}\\
    \phantom{*} & \phantom{*} & |[uband]| * & |[uband]| * & \phantom{*}\\
    \phantom{*} & \phantom{*} & \phantom{*} & |[uband]| * & |[uband]| *\\
    \phantom{*} & \phantom{*} & \phantom{*} & \phantom{*} & |[uband]| *\\
  };
  \node[title, above=5pt] at (U.north) {$U$};
  \end{scope}
\end{tikzpicture}

  \caption{A tridiagonal matrix can factor as bidiagonal $L$ times bidiagonal $U$, retaining local band structure.}
  \label{fig:band-matrix-lu-structure}
\end{figure}

\begin{definition} \label{def:band-and-tridiagonal-matrix} \hlt{Band and Tridiagonal Matrices}\\
For $A=[a_{ij}]\in\operatorname{Mat}(n,n,\mathbb{F})$ and $1\leq w\leq n$, say that $A$ has \hlt{half-bandwidth} $w$ if $a_{ij}=0$ whenever $\abs{i-j}\geq w$. Thus $w=1$ gives a diagonal matrix, while $w=2$ gives a \hlt{tridiagonal} matrix.
\end{definition}

\begin{method} \label{meth:banded-elimination-preview} \hlt{Banded Elimination Preview}\\
Let $A\in\operatorname{Mat}(n,n,\mathbb{F})$ have half-bandwidth $w$.
\begin{enumerate}[label=(\roman*)]
\setlength{\itemsep}{0pt}
\item Eliminate without row exchanges that leave the band.
\item At pivot $k$, update only entries with $\abs{i-k}<w$ and $\abs{j-k}<w$.
\item Obtain $A=LU$ with $L$ and $U$ of comparable bandwidth.
\end{enumerate}
Cost: $O(nw^2)$ operations; for tridiagonal matrices $(w=2)$, $O(n)$.
\end{method}

\begin{perspective} \label{per:five-point-stencil-bandwidth} \hlt{Five-Point Stencils Are Banded but Less Local}\\
The $1$D second-difference pattern $-1,2,-1$ is tridiagonal; the row-ordered $N\times N$ five-point stencil is still sparse, but vertical neighbors create diagonals offset by about $N$. Thus bandwidth grows with grid width, and incomplete LU keeps a prescribed sparse pattern for a cheap approximate solve or preconditioner.
\end{perspective}

\begin{figure}[H]
  \centering
  \begin{tikzpicture}[
  >=Stealth,
  every node/.style={font=\small},
  box/.style={draw=gray!50, fill=gray!3, rounded corners=2pt, line width=0.55pt},
  main/.style={line width=1.0pt, black},
  near/.style={line width=0.8pt, blue!70!black},
  far/.style={line width=0.8pt, teal!70!black},
  bracket/.style={gray!75!black, line width=0.55pt},
  legend/.style={font=\footnotesize, anchor=west}
]
  \node at (-0.62,0) {$A=$};

  \draw[bracket] (-0.18,1.32) -- (-0.34,1.32) -- (-0.34,-1.32) -- (-0.18,-1.32);
  \draw[bracket] (3.18,1.32) -- (3.34,1.32) -- (3.34,-1.32) -- (3.18,-1.32);

  \draw[box] (0,-1.2) rectangle (3,1.2);
  \begin{scope}
    \clip (0,-1.2) rectangle (3,1.2);
    \draw[far] (-0.1,0.45) -- (1.75,-1.35);
    \draw[near] (-0.1,0.95) -- (2.45,-1.35);
    \draw[main] (0.25,1.25) -- (2.9,-1.35);
    \draw[near] (0.75,1.35) -- (3.15,-0.95);
    \draw[far] (1.55,1.35) -- (3.15,-0.20);
  \end{scope}

  \draw[main] (3.9,0.72) -- (4.35,0.72);
  \node[legend] at (4.45,0.72) {offset $0$};
  \draw[near] (3.9,0.12) -- (4.35,0.12);
  \node[legend, blue!70!black] at (4.45,0.12) {offset $\pm1$};
  \draw[far] (3.9,-0.48) -- (4.35,-0.48);
  \node[legend, teal!70!black] at (4.45,-0.48) {offset $\pm N$};
\end{tikzpicture}

  \caption{Row ordering gives matrix bands at offsets $0,\pm1,\pm N$.}
  \label{fig:five-point-stencil-bandwidth}
\end{figure}
